
\section{用椭圆曲线构造的反例与病态}

现在，我们要利用对\trEllipticCurve{}的理解来给出一些有趣的反例。\trEllipticCurve{}能提供给我们的额外汁水，来自这样一个事实：\trEllipticCurve{}是具有「连续\trPicardGroup」的最简单的\trVariety。

\textbf{19.11.1. 一个\trFactorial{}的概形，却没有哪一点的仿射开邻域是某个\trUFD{}的 $\Spec$。}

设 $E$ 是 $\mathbb{C}$ 上（或某个其他不可数代数闭域上的）一条\trEllipticCurve。取 $p\in E$。\trLocalRing{} $\OO_{E,p}$ 是一个\trDVR，从而是一个\trUFD（\mbox{定理~13.5.8}）。那么 $E$ 的任何仿射开邻域都形如 $E\setminus q_1\setminus\cdots\setminus q_n$，其中 $q_1$、$\dots$、$q_n$ 为一些点。我断言：它的\trPicardGroup{}是非平凡的。回想那条正合列：
\[
\mathbb{Z}^{\oplus n}\mathrel{\xrightarrow{\ (a_1,\dots,a_n)\,\mapsto\, a_1q_1+\cdots+a_nq_n\ }}\mathrm{Pic}\,E\longrightarrow\mathrm{Pic}(E\setminus q_1\setminus\cdots\setminus q_n)\longrightarrow 0.
\]
但左边的群是可数的，而中间的群是不可数的，于是右边的群，即 $\mathrm{Pic}(E\setminus q_1\setminus\cdots\setminus q_n)$，是非零的。我们已经证明：$E$ 的每个非空\trOpenSubset{}都有非零的\trLineBundle，一如\mbox{注记~14.2}所许诺\footnote{原书此处作「Remark 14.2」，而 \S14.2 中并无编号为「14.2」的注记；这一许诺实出于 \S14.2 的引言段（p.405）。兹照录原书。——译者注}。

由于 $E\setminus q_1\setminus\cdots\setminus q_n$ 是\trAffine{}的，据\mbox{练习~15.5.P}，与之相应的环不是一个\trUFD。总结一下：复\trEllipticCurve{}是\trFactorial{}的，却没有一个仿射\trOpenSubset{}带有一个唯一分解的环，一如 \S5.4.4 所许诺。

\textbf{19.11.A.（练习）}\quad 上面的论证表明：在不可数域上，$\mathrm{Pic}\,E$ 不是有限生成的\trGroup。请证明：即使在可数域 $\bar{\mathbb{Q}}$ 上，$\mathrm{Pic}\,E$ 也不是有限生成的，思路如下。若\trEllipticCurve{} $E$ 由 $q_1$、$\dots$、$q_n$ 生成，则存在 $\mathbb{Q}$ 的一个有限域扩张 $K$，使得诸 $q_i$ 全都定义于其上——它就是诸 $q_i$ 的\trResidueField{}的\trCompositum。（当心——你明白诸 $q_i$「定义于」一个比 $\bar{\mathbb{Q}}$ 更小的域之上是什么意思吗？）证明：$E$ 的由 $q_1$、$\dots$、$q_n$ 生成的\trSubgroup{}中的任何点也必定定义于 $K$ 之上。证明：$E$ 有一个不定义于 $K$ 之上的点。利用这些来证明：$\mathrm{Pic}\,E$ 不是有限生成的。（把 $\bar{\mathbb{Q}}$ 换成 $\bar{\mathbb{F}}_p$，同样的论证也适用。）

\textbf{19.11.2.}\ \emph{\trhRemark{}。}\quad 与上面的讨论相对照：在 $\mathbb{Q}$ 上，\trn{Mordell--Weil}{莫德尔--韦伊}定理断言 $\mathrm{Pic}\,E$ 是有限生成的（见\mbox{旁注~19.9.4}）。

\textbf{19.11.3.}\ \extrahard\ \textbf{一个带无穷多种代数结构的复曲面（供具备足够复几何背景者）。}

一如 \S7.3.16 中所述，一个\trComplex\trManifold{}可以带有许多种代数结构。下面这个出自 M.~Kim 的例子给出了一个带\emph{无穷}多种代数结构的例子。设 $E$ 是 $\mathbb{C}$ 上一条带有原点 $p$ 的\trEllipticCurve，并设 $\mathscr{L}$ 是 $E-p$ 上的一个非平凡\trLineBundle。那么 $\mathscr{L}$ 是\emph{解析}平凡的，因为 $E-p$ 是一个\trStein{}空间，故这个\trTotalSpace{}的\trAnalytification{}不依赖于 $\mathscr{L}$ 的选取。然而，你所掌握的知识已足以让你（费一番功夫）证明：存在无穷多个两两（在代数上）互不同构的 $\mathscr{L}$，而且它们的全空间也两两（在代数上）互不同构。这便给出了一个带无穷多种代数结构的\trComplex\trSurface（事实上其数目有连续统之多）。详情见 [MO68421]。

\textbf{19.11.4. 用非挠点构造的反例。}

接下来，我们要利用\trComplex\trEllipticCurve{}的非挠点的存在性，给出若干反例。我们先来证明这样一个点的存在性。

\textbf{19.11.5.}\quad 我们有「乘以 $n$」映射 $[n]\colon E\to E$，它把 $p$ 送到 $np$。若 $n=0$，则映射 $[n]$ 的\trDegree{}为 $0$（即使 $[n]$ 并非\trDominant{}的，\trDegree{}仍有定义；见\mbox{定义~12.2.2}）。若 $n=1$，则映射 $[n]$ 的\trDegree{}为 $1$。鉴于环面的复图景，你对 $[n]$ 的\trDegree{}是 $n^2$ 或许不会感到惊讶。若 $n=2$，我们差不多已经证明了它的\trDegree{}为 $4$，因为我们已经验证过：恰有四个点 $q$ 使得 $2p=2q$（\mbox{练习~19.9.A}）。而这真正表明的（利用\mbox{练习~16.3.D(b)}）不过是\trDegree{}至少为 $4$。（假如我们真想这么做的话，本可以动手验证\trDegree{}确实是 $4$\footnote{原书此处作「the degree is 4 is we really wanted to」，「is」当为「if」之误，径改。——译者注}。）

\textbf{19.11.6.\ \trhProposition。}\ ——\ \emph{设 $E$ 是特征非 $2$ 的域 $k$ 上的一条\trEllipticCurve。对每个 $n>0$，「乘以 $n$」\trMorphism{} $[n]$ 都有正\trDegree，于是 $n$-挠点只有有限多个。}

\emph{证明。}\quad 我们不妨设 $k=\bar{k}$，因为曲线之间映射的\trDegree{}不依赖于域扩张。

我们对这个结果作归纳证明；它在 $n=1$ 与 $n=2$ 时已然成立。

若 $n$ 是奇数（比如说 $2k+1$），则反设：对一切\trClosedPoint{} $q$ 都有 $nq=0$。取 $r$ 为一个非平凡的 $2$-挠点，于是 $2r=0$。但 $nr=0$ 同样成立，于是 $r=(n-2k)r=0$，与 $r\ne 0$ 矛盾。

若 $n$ 是偶数，则 $[n]=[2]\circ[n/2]$（\trDegree{}在\trRationalMap{}的\trComposition{}下是乘性的，\S12.2.2），而由归纳假设，$[2]$ 与 $[n/2]$ 都有正\trDegree。\hfill$\square$

特别地，$E$ 上挠点的总数是可数的。若 $k$ 是不可数域，则 $E$ 有不可数多个\trClosedPoint（考虑该\trCurve{}的形如 $y^2=x^3+ax+b$ 的一个开子集；$x$ 的取法有不可数多种，而每一种取法又有一到两种 $y$ 的取法）。于是我们便有如下结论。

\textbf{19.11.7.\ \trhCorollary。}\ ——\ \emph{若 $E$ 是不可数代数闭域（例如 $\mathbb{C}$）上特征非 $2$ 的一条\trCurve，则 $E$ 有一个非挠点。}

\emph{证明。}\quad 对每个 $n$，$n$-挠点只有有限多个。于是挠点（至多）可数。而\trCurve{} $E$ 有不可数多个\trClosedPoint。（一个论证：取一个二重覆盖 $\pi\colon E\to\mathbb{P}^1$。那么 $\mathbb{P}^1$ 有不可数多个\trClosedPoint，而 $\pi$ 在\trClosedPoint{}上是满的。）\hfill$\square$

\textbf{19.11.8.}\ \emph{\trhRemark{}。}\quad 这个论证在可数域上显然失效。事实上，在 $\mathbb{F}_p$ 上，\trEllipticCurve{} $E$ 的所有点都是挠点。（任何一个给定的点都定义于某个有限域 $\mathbb{F}_{p^r}$ 之上。定义于 $\mathbb{F}_{p^r}$ 之上的点构成 $E$ 的一个\trSubgroup——利用群律的显式几何构造可知——而在 $\mathbb{F}_{p^r}$ 之上的点只有有限多个\footnote{原书此处作「there are finite number of points」，当为「there are a finite number of points」之误（脱冠词），径改。——译者注}——当然不会多于 $\mathbb{P}^2$ 的 $\mathbb{F}_{p^r}$-点的个数。）但在 $\bar{\mathbb{Q}}$ 上，确有带非挠点的\trEllipticCurve。更好的是，在 $\mathbb{Q}$ 上就有例子：$[2,1,8]$ 是\trEllipticCurve{} $y^2z=x^3+4xz^2-z^3$ 的一个 $\mathbb{Q}$-点，它不是挠点。这个证明会把我们带得太远，但有一种方法是使用 Nagell--Lutz 定理（例如见 [Si1, Cor.~7.2]）。

现在，我们利用非挠点的存在性来制造一些有趣的病态。

\textbf{19.11.9. 一个并非由分次环映射产生的射影簇态射，即便\trRegrading{}之后亦然。}

\textbf{19.11.B.（练习）}\quad 设 $E\hookrightarrow\mathbb{P}^2_{\mathbb{C}}$ 是一条光滑的\trComplex{}平面三次\trCurve（从而\trGenus{}为 $1$），随之便得到一个\trGradedRing{} $S_\bullet$。设 $t\colon E\to E$ 是由一个非挠点给出的\trTranslation。证明：$t$ 不对应于任何\trGradedRingMap{} $S_\bullet\to S_\bullet$，即便\trRegrading{}之后也不对应（对照\mbox{注记~15.2.9}）。（如果你愿意，还可以证明下面这件事。若 $u\colon E\to E$ 是一个\trTranslation，证明：$u$ 对应于某个重新分次的环之间的映射 $S_{n\bullet}\to S_{n\bullet}$，当且仅当 $u$ 是由一个挠点给出的\trTranslation。）

\textbf{19.11.10. 一个\trAffine\trScheme{}的不是特异开集的仿射\trOpenSubset。}

我们可以用这一点来构造一个例子：一个\trAffine\trScheme{} $X$ 连同一个在 $X$ 中不是特异开集的仿射\trOpenSubset{} $Y$。令 $X=E-p$，它是\trAffine{}的（见\mbox{练习~19.2.C}；或者更好的做法是，注意线性系 $\OO(3p)$ 把 $E$ 送到 $\mathbb{P}^2$ 的方式使得「无穷远线」只与 $E$ 相交于 $p$；于是 $E-p$ 便有一个到仿射概形 $\mathbb{A}^2$ 的\trClosedEmbedding）。

设 $q$ 是 $E$ 上的另一点，使得 $q-p$ 非挠。那么 $E-p-q$ 是\trAffine{}的（\mbox{练习~19.2.C}）。假设它是特异开集。那么 $E-p$ 上存在一个函数 $f$，它在 $q$ 处消失（以某个正的阶 $d$）。于是 $f$ 是 $E$ 上的一个\trRationalFunction，它在 $q$ 处以 $d$ 阶消失，并且（由于 $f$ 的\trZeros{}总数减去\trPole{}总数为 $0$）在 $p$ 处有一个 $d$ 阶\trPole。但这样一来，在 $\mathrm{Pic}^0E$ 中就有 $d(p-q)=0$，与我们关于 $p-q$ 非挠的假设相矛盾。

特别地，$E-p$ 是一个\trAffine\trScheme，而 $q$ 局部地由单个方程切出，但它\emph{即便在集合论的层面上}也不是由单个方程整体切出的。这一点曾在\mbox{练习~8.3.E}与 \S9.5.1 中许诺。

\textbf{19.11.11. 一个没有非平凡\trLineBundle{}的\trProper（非射影）\trSurface。}

接下来，我们利用\trEllipticCurve{} $E$ 上的一个非挠点 $p$ 来构造一个没有非平凡\trLineBundle{}的\trProper{}非射影\trSurface。令 $X_1$ 为 $\mathbb{P}^2$，$Z_1$ 为 $X_1$ 中的一条光滑三次\trCurve，将它等同于 $E$。令 $Z_2\subset X_2$ 完全同样。沿着同构 $Z_1\overset{\sim}{\longleftrightarrow}Z_2$ 把 $X_1$ 与 $X_2$ 粘合起来，这个同构不由恒等给出，而由 $p$ 的\trTranslation{}给出。（这一粘合构造曾在 \S15.2.12 中描述。）把所得的结果称作 $X$。据\mbox{练习~15.2.M}（及其后那句话），$X$ 是\trProper{}的。

\textbf{19.11.C.（练习）}\quad 证明：$X$ 上的每条\trLineBundle{}都是平凡的。提示：设 $\mathscr{L}$ 是 $X$ 上的一个\trInvertibleSheaf。那么 $\mathscr{L}|_{X_i}$ 是某个 $\OO(d_i)$（$i=1$、$2$）。$\mathscr{L}|_{X_1}$ 到 $Z_1$ 上的限制必须与 $\mathscr{L}|_{X_2}$ 到 $Z_2$ 上的限制相一致。利用这一点证明：$d_1=d_2=0$。并解释：为什么以这种方式把两个平凡的丛（分别在 $X_1$ 与 $X_2$ 上）粘合起来，得到的便是 $X$ 上的一个平凡丛。

\textbf{19.11.D.（练习）}\quad 证明：$X$ 不是射影的。

一个稍微简单些的\trProper{}非射影\trSurface{}的例子，见\mbox{练习~20.2.G}。

\textbf{19.11.12. 一个绝无可能成为概形的\trPicardGroup。}

我们曾非正式地观察到：一条\trEllipticCurve{}的\trPicardGroup{}「打心底里想成为一个概形」（见 \S19.10.5）。对一般的射影（甚至是\trProper{}的）\trVariety{}来说，这确实成立（见 [FGIKNV, Ch.~9]）。另一方面，若我们在 $\mathbb{C}$ 上工作，仿射概形 $E-p-q$（用上面 \S19.11.10 的语言来说）的\trPicardGroup{}可以被解读为 $\mathbb{C}$ 模掉一个\trLattice{}再模掉一个非挠点（例如 $\mathbb{C}/\langle 1,i,\pi\rangle$）。作为流形，它没有任何合理的解读，更遑论作为\trVariety{}了。所以，\trProper{}\trVariety{}的\trPicardGroup{}居然当真是一个概形，这一事实应当被看作十分了不起的。

\textbf{19.11.13. 一个其函数环不是有限生成的\trVariety。}

接下来我们给出一个\trComplex\trVariety{}的例子，其\trGlobalSections{}环不是有限生成的。（在 $\mathbb{Q}$ 上的例子可以用同样的方式、利用\mbox{注记~19.11.8} 中的那条曲线构造出来。如果你愿意，你还可以证明这些例子甚至是\trQuasiaffine{}的。）这与\trHilbert{}第十四问题有关（见 [Mu5, \S3]）。

\textbf{19.11.E.（预备练习）}\quad 设 $X$ 是一个概形，而 $L$ 是对应于\trInvertibleSheaf{} $\mathscr{L}$ 的那条\trLineBundle{}的\trTotalSpace，于是 $L=\mathit{Spec}\oplus_{n\ge0}(\mathscr{L}^{\vee})^{\otimes n}$。（这个构造最初出现于\mbox{定义~17.1.5}。）证明：若 $X$ 是\trQuasicompact{}的，则 $\mathrm{H}^0(L,\OO_L)=\oplus\mathrm{H}^0(X,(\mathscr{L}^{\vee})^{\otimes n})$。（可能的提示：为 $\mathscr{L}$ 选取一个\trTrivialization{}覆盖。反问一句：如果 $\mathscr{L}$ 是\trRank{}为 $r$ 的\trLocallyFreeSheaf，你能想出更一般的表述是什么吗？）

设 $E$ 是某个\trBaseField{} $k$ 上的一条\trEllipticCurve，$\mathscr{N}$ 是 $E$ 上\trDegree{}为 $0$ 的一个非挠\trInvertibleSheaf，而 $\mathscr{P}$ 是 $E$ 上一个正\trDegree{}的\trInvertibleSheaf。那么 $\mathrm{H}^0(E,\mathscr{N}^m\otimes\mathscr{P}^n)$ 非零，当且仅当下面二者之一成立：(i) $n>0$；或者 (ii) $m=n=0$（此时诸\trSection{}都是 $k$ 中的元素）。

\textbf{19.11.F.（容易的练习）}\quad 证明：环 $R=\oplus_{m,n\ge0}\mathrm{H}^0(E,\mathscr{N}^m\otimes\mathscr{P}^n)$ 不是有限生成的。

\textbf{19.11.G.（练习）}\quad 设 $X$ 是 $E$ 上相伴于 $(\mathscr{N}\oplus\mathscr{P})^{\vee}$ 的\trVectorBundle{}的\trTotalSpace。证明：$X$ 的\trGlobalSections{}环就是 $R$，从而不是有限生成的。（提示：把 $X$ 解读为 $E$ 上一条\trLineBundle{}之上的\trLineBundle。）

\textbf{19.11.H.（练习）}\quad 证明：$X$（如上一道练习中那样）是一个\trVariety（并且特别地是\trNoetherian{}的），其\trGlobalSections{}环不是\trNoetherian{}环。
